% Claim statement file for row claim_3_6 -- "SplitTopologicalOrder".
% Auto-generated stub by scaffold/solve_chapter.py at row start. The
% manager agent (or a worker it dispatches) fills the body below with the
% claim's statement(s) from the lecture notes -- statement only, no proof
% (proofs live in the sibling `claim_3_6_proof_*.tex` file).
%
% This file is a sibling of `main.tex` inside `<subsection>/tex/`. The
% `subfiles` package lets it render both standalone and as part of
% `main.tex`. Note: `main.tex` does NOT \subfile this statement file -- the
% sibling `_proof_` file restates the statement above the proof, so
% including both in `main.tex` would be redundant. This file exists for
% standalone reading / grep convenience.

\documentclass[main]{subfiles}

\begin{document}
% ref: claim_3_6 (statement)
% title: SplitTopologicalOrder
\def\rowref{claim\_3\_6}
\phantomsection\label{claim_3_6}

\begin{Rem}[SplitTopologicalOrder]
    Let $G = (J, V, E, L)$ be a CADMG, i.e.\ a CDMG (def \ref{def-cdmg}) that is acyclic in the sense of def \ref{def-acylic} (equivalently, $G$ satisfies the attribute ``CADMG'' of def \ref{def-cdmg-types}, item~i.); throughout we use the notation of def \ref{not-cdmg}, the family-relationship notation of def \ref{def_3_5}, the topological-order definition of def \ref{def-topological-order}, and the node-splitting construction of def \ref{def:G_node-splitting}. Let
    \[ W \ins V \]
    be a subset of the output-node set of $G$, and let
    \[ G_{\spl(W)} \;=\; \lp J_{\spl(W)}, V_{\spl(W)}, E_{\spl(W)}, L_{\spl(W)} \rp \]
    be the node-split graph of $G$ w.r.t.\ $W$ as constructed in def \ref{def:G_node-splitting}. Both quantifiers ``for every CADMG $G$'' and ``for every $W \ins V$'' are read \emph{universally}, with the side condition $W \ins V$ inherited verbatim from def \ref{def:G_node-splitting}; the case $W = \emptyset$ is admitted (then $W^0 = W^1 = \emptyset$ and $G_{\spl(W)} = G$ by items i.--iv.\ of def \ref{def:G_node-splitting}).

    \emph{Carrier of $G_{\spl(W)}$.} By items i.\ and ii.\ of def \ref{def:G_node-splitting},
    \[
        J_{\spl(W)} \cup V_{\spl(W)} \;=\; J \cup (V \sm W) \dcup W^0 \dcup W^1,
    \]
    where $W^0 := \lC w^0 \st w \in W \rC$ and $W^1 := \lC w^1 \st w \in W \rC$ are the two tagged copies of $W$ introduced in def \ref{def:G_node-splitting}. The four constituents $J$, $V \sm W$, $W^0$, $W^1$ are pairwise type-level disjoint by def \ref{def:G_node-splitting}: $J \cap V = \emptyset$ (def \ref{def-cdmg}), $(V \sm W) \cap W = \emptyset$, and the tagged copies $W^0$, $W^1$ are formally distinct from each element of $J \cup V$ and from each other ($W^0 \cap W^1 = \emptyset$, $w^0 \neq w^1$ for every $w \in W$).

    \emph{Finiteness and indexing of $J \cup V$.} By def \ref{def-cdmg} items~i.--ii., the set $J \cup V$ is finite; set
    \[ n \;:=\; |J \cup V|. \]

    \emph{The remark.} The following conjunction of sub-claims holds.

    \begin{enumerate}[label=(\alph*)]
        \item \emph{$G_{\spl(W)}$ is acyclic.}
              The CDMG $G_{\spl(W)}$ is acyclic in the sense of def \ref{def-acylic}, i.e.\ for every $x \in J_{\spl(W)} \cup V_{\spl(W)} = J \cup (V \sm W) \dcup W^0 \dcup W^1$, there does not exist any non-trivial directed walk from $x$ to itself in $G_{\spl(W)}$.

        \item \emph{An explicit topological order on $G_{\spl(W)}$ obtained from any topological order on $G$.}
              For every topological order $<$ of $G$ in the sense of def \ref{def-topological-order}, presented in the equivalent indexing form (def \ref{def-topological-order}, ``Equivalent indexing form'') as a bijection
              \[
                  \{1, 2, \dots, n\} \;\longrightarrow\; J \cup V, \qquad i \;\longmapsto\; v_i,
              \]
              with $v_1 < v_2 < \cdots < v_n$ (so that ``parents precede children'' holds in this enumeration: $v_i \in \Pa^G(v_j) \Rightarrow i < j$), the binary relation $<_{\spl}$ on $J_{\spl(W)} \cup V_{\spl(W)}$ defined below is a topological order of $G_{\spl(W)}$ in the sense of def \ref{def-topological-order}.

              \emph{Index map.}
              Define the index map
              \[ \phi \,:\, J_{\spl(W)} \cup V_{\spl(W)} \;\longrightarrow\; \Q \]
              by case analysis on the disjoint-union carrier $J \cup (V \sm W) \dcup W^0 \dcup W^1$:
              \begin{itemize}
                  \item \emph{Unsplit nodes retain their original index.}
                        For every $x \in J \cup (V \sm W)$, write $x = v_j$ with $j \in \{1, 2, \dots, n\}$ the unique index given by the enumeration of $J \cup V$ above, and put
                        \[ \phi(x) \;:=\; j. \]
                  \item \emph{Tagged copies of $W$ are placed symmetrically around the original index.}
                        For every $w \in W$, write $w = v_j$ with $j \in \{1, 2, \dots, n\}$ the unique index given by the enumeration of $J \cup V$ above, and put
                        \[
                            \phi(w^0) \;:=\; j - \tfrac{1}{3}, \qquad \phi(w^1) \;:=\; j + \tfrac{1}{3}.
                        \]
              \end{itemize}

              \emph{Order from index map.}
              Define
              \[
                  x \;<_{\spl}\; y \quad :\Longleftrightarrow\quad \phi(x) \,<\, \phi(y) \qquad \text{for } x, y \in J_{\spl(W)} \cup V_{\spl(W)},
              \]
              where the right-hand inequality is the standard strict order on $\Q$.
    \end{enumerate}

    \emph{Where each clause of the LN's prose lands.}
    The LN's clauses ``$v_j^0$ the index $j - \tfrac{1}{3}$'' and ``$v_j^1$ the index $j + \tfrac{1}{3}$'' for $v_j \in W$ are the second bullet of the index map above. The LN's closing prose ``and then ordering all nodes according to their index value'' is the definition of $<_{\spl}$ from $\phi$ via the strict order on $\Q$. The first bullet of the index map -- ``unsplit nodes $v_j \in J \cup (V \sm W)$ retain index $j$'' -- is the rule the LN leaves implicit: the LN's literal text assigns indices only to the tagged copies $v_j^0$, $v_j^1$ of split nodes $v_j \in W$ and says nothing about the indices of $v_j \in J \cup (V \sm W)$, but the closing instruction ``order all nodes according to their index value'' is well-defined on the carrier $J_{\spl(W)} \cup V_{\spl(W)}$ only if the unsplit nodes also carry an index. The natural and only-consistent reading is that the original index $j$ is retained for unsplit nodes; the first bullet records this retention rule explicitly.

    \emph{Compatibility of the tagging convention with def \ref{def:G_node-splitting}.}
    For every $w \in W$ with $w = v_j$, $\phi(w^0) = j - \tfrac{1}{3} < j + \tfrac{1}{3} = \phi(w^1)$, so $w^0 <_{\spl} w^1$. This is consistent with item iii.\ of def \ref{def:G_node-splitting}: that item adds the \emph{transfer} edge $(w^0, w^1) \in E_{\spl(W)}$ for every $w \in W$, i.e.\ $w^0 \in \Pa^{G_{\spl(W)}}(w^1)$ in the sense of def \ref{def_3_5}, item~i., and the parent-precedence clause of def \ref{def-topological-order} for $G_{\spl(W)}$ then demands precisely $w^0 <_{\spl} w^1$. (Under the opposite tagging convention -- transfer edge $(w^1, w^0)$, incoming edges into $w$ reattached at $w^1$, outgoing edges out of $w$ leaving $w^0$ -- the symmetric index assignment ``$j - \tfrac{1}{3}$ for the $^0$-copy and $j + \tfrac{1}{3}$ for the $^1$-copy'' would not produce a topological order; sub-claim~(b) is correct as stated precisely because the index choice matches def \ref{def:G_node-splitting}'s tagging convention.)

    \emph{Role of the CADMG (acyclicity) precondition.}
    The construction of def \ref{def:G_node-splitting} on its own does not preserve acyclicity from $G$ to $G_{\spl(W)}$: if $(w, w) \in E$ for some $w \in W$, then the first set-builder in item iii.\ of def \ref{def:G_node-splitting} adds the lifted directed edge $(w^1, w^0) \in E_{\spl(W)}$, which together with the transfer edge $(w^0, w^1) \in E_{\spl(W)}$ realises a length-$2$ directed cycle $w^0 \to w^1 \to w^0$ in $G_{\spl(W)}$ (cf.\ the closing paragraph ``Self-loops on $W$ produce $2$-cycles, and $G_{\spl(W)}$ is still a CDMG'' of def \ref{def:G_node-splitting}). The CADMG hypothesis on $G$ -- i.e.\ acyclicity of $G$ in the sense of def \ref{def-acylic} -- rules this hazard out: by the consequence ``no directed self-loops on output nodes'' established in def \ref{def-acylic}, an acyclic CDMG has $(v, v) \notin E$ for every $v \in V$, and in particular $(w, w) \notin E$ for every $w \in W \ins V$. Thus the directed-edge set $E_{\spl(W)}$ is free of the $2$-cycle hazard of def \ref{def:G_node-splitting}, which is what allows sub-claim~(a) to hold. The bidirected-edge set $L$ plays no role in either sub-claim: acyclicity (def \ref{def-acylic}) is defined via the non-existence of non-trivial directed walks (which use only $E$), and a topological order (def \ref{def-topological-order}) is defined via the parent-set $\Pa^G$ (def \ref{def_3_5}, item~i., which is defined from $E$ alone).
\end{Rem}

\end{document}
